\chapter{绪论}
\label{chap:intro}

\section{研究背景与意义}
城市交通流量预测是交通管理的基础任务。示例城市的检测器覆盖率不足，观测数据稀疏。

\begin{figure}[htbp]
  \centering
  \includegraphics[width=0.7\textwidth]{chapter1/background.png}
  \bicaption{示例城市检测器分布}{Distribution of detectors in the example city}
  \label{fig:c1-background}
\end{figure}

\section{国内外研究现状}
\begin{figure}[htbp]
  \centering
  \includegraphics[width=0.7\textwidth]{chapter1/literature.png}
  \bicaption{示例文献数量统计}{Statistics of example publications}
  \label{fig:c1-literature}
\end{figure}

\subsection{基于统计模型的预测方法}
统计模型结构简单，但难以描述非线性关系。

\subsection{基于深度学习的预测方法}
深度学习方法精度较高，但依赖完整观测。

\subsection{稀疏观测下的预测方法}
现有稀疏观测方法多针对单一数据源。% 这是一条注释

\section{研究问题与关键挑战}
\begin{enumerate}
  \item 观测稀疏导致空间相关性难以刻画。
  \item 单一数据源难以覆盖复杂场景。
  \item 模型部署后缺少在线更新机制。
\end{enumerate}

\section{研究内容与章节安排}
\begin{figure}[htbp]
  \centering
  \includegraphics[width=0.8\textwidth]{chapter1/organization.png}
  \bicaption{论文组织结构}{Organization of the thesis}
  \label{fig:c1-organization}
\end{figure}
